% ─────────────────────────────────────────────────────────────────────────────
% SHARED coursework preamble — every course inputs THIS.
%
% PREAMBLE ONLY. No \begin{document}, no course-specific notation.
%
% Course-specific macros live at COURSE level, in
%   content/<course>/reference_docs/<course>_macros.tex
% which is inputted AFTER this file. ONE per course, shared by hw/qz/exam --
% macros under a single kind are unreachable from a sibling kind.
%
% That split is the point: the page style, the answer box, and the problem
% header are identical everywhere, so they are defined once; \conv and \ustep
% mean something in DSP and nothing in a networks course, so they are not.
%
% A per-problem file therefore starts:
%   % ─────────────────────────────────────────────────────────────────────────────
% SHARED coursework preamble — every course inputs THIS.
%
% PREAMBLE ONLY. No \begin{document}, no course-specific notation.
%
% Course-specific macros live at COURSE level, in
%   content/<course>/reference_docs/<course>_macros.tex
% which is inputted AFTER this file. ONE per course, shared by hw/qz/exam --
% macros under a single kind are unreachable from a sibling kind.
%
% That split is the point: the page style, the answer box, and the problem
% header are identical everywhere, so they are defined once; \conv and \ustep
% mean something in DSP and nothing in a networks course, so they are not.
%
% A per-problem file therefore starts:
%   % ─────────────────────────────────────────────────────────────────────────────
% SHARED coursework preamble — every course inputs THIS.
%
% PREAMBLE ONLY. No \begin{document}, no course-specific notation.
%
% Course-specific macros live at COURSE level, in
%   content/<course>/reference_docs/<course>_macros.tex
% which is inputted AFTER this file. ONE per course, shared by hw/qz/exam --
% macros under a single kind are unreachable from a sibling kind.
%
% That split is the point: the page style, the answer box, and the problem
% header are identical everywhere, so they are defined once; \conv and \ustep
% mean something in DSP and nothing in a networks course, so they are not.
%
% A per-problem file therefore starts:
%   \input{../../../../docs/latex/coursework_preamble.tex}
%   \input{../../reference_docs/cesc470_macros.tex}
%   \begin{document}
%
% Build-check one file, or a whole assignment:
%   docs/latex/build_tex.sh content/cesc_470/hw/hw01/p01_five_components.tex
%   docs/latex/build_tex.sh content/cesc_470/hw/hw01
%
% Doctrine: docs/directives/coursework-solutions.md
% ─────────────────────────────────────────────────────────────────────────────

\documentclass[11pt]{article}

\usepackage[margin=1in]{geometry}
\usepackage{amsmath, amssymb, mathtools}
\usepackage{graphicx}
\usepackage{float}
\usepackage{booktabs}
\usepackage{longtable}
\usepackage{enumitem}
\usepackage{siunitx}
\usepackage[dvipsnames]{xcolor}
\usepackage{tcolorbox}
\usepackage{fancyhdr}
\usepackage{caption}
\usepackage{tikz}
\usepackage{pgfplots}
\pgfplotsset{compat=1.18}
\usepackage[colorlinks=true, allcolors=NavyBlue]{hyperref}

\captionsetup{font=small, labelfont=bf}
\setlist[enumerate]{itemsep=2pt, topsep=4pt}
\setlist[itemize]{itemsep=2pt, topsep=4pt}

% --- final-answer box -------------------------------------------------------
% Every problem file ends with its result in one of these. The condensed
% solutions document is assembled by lifting exactly these.
\newtcolorbox{answerbox}[1][]{
  colback=Green!6, colframe=Green!55!black,
  boxrule=0.6pt, arc=2pt, left=6pt, right=6pt, top=4pt, bottom=4pt,
  #1
}

% --- citation to course material --------------------------------------------
% \cite{Module 01, PDF p55}  ->  a small italic parenthetical.
% Solutions cite the SLIDE/PAGE a formula came from so a grader (or a future
% reader) can check the claim against what was actually taught. See the
% directive: every non-obvious step carries one.
\newcommand{\srcref}[1]{{\small\itshape(#1)}}

% --- a step that came from OUTSIDE the course materials ---------------------
% Used only where the course is genuinely silent and the problem asks for
% outside research. Marked visually so it is never mistaken for lecture content.
\newcommand{\extref}[1]{{\small\itshape\color{Sepia}[external: #1]}}

% --- callout for the trap in a problem --------------------------------------
\newtcolorbox{trap}[1][]{
  colback=BurntOrange!7, colframe=BurntOrange!70!black,
  boxrule=0.6pt, arc=2pt, left=6pt, right=6pt, top=4pt, bottom=4pt,
  fonttitle=\bfseries, title=Watch out, #1
}

% --- per-problem header -----------------------------------------------------
% \problemheader{8}{15 pts}{Target clock rate}
% The optional 4-argument form carries a learning outcome where the course
% states one (CESC 410 does; CESC 470's HW1 does not).
\newcommand{\problemheader}[3]{%
  \begin{center}\large\bfseries
    Problem #1 \quad\normalsize\mdseries(#2)\\[2pt]
    \large\bfseries #3
  \end{center}
  \vspace{2pt}\hrule\vspace{8pt}
}
\newcommand{\problemheaderlo}[4]{%
  \begin{center}\large\bfseries
    Problem #1 \quad\normalsize\mdseries(#2, #3)\\[2pt]
    \large\bfseries #4
  \end{center}
  \vspace{2pt}\hrule\vspace{8pt}
}

% Course name in the running head. Defined BEFORE \lhead references it, and
% \renewcommand'd by each course's macros file (inputted after this one).
\providecommand{\coursename}{}

\pagestyle{fancy}
\fancyhf{}
\lhead{\small\coursename}
\rhead{\small\thepage}
\renewcommand{\headrulewidth}{0.3pt}

%   % CESC 470 Computer Architecture — course-specific macros ONLY.
%
% Inputted AFTER docs/latex/coursework_preamble.tex, which carries everything
% shared (answerbox, problemheader, srcref, page style). Put nothing general
% here: if another course would want it, it belongs in the shared preamble.

\renewcommand{\coursename}{CESC 470}

% --- performance-equation notation ------------------------------------------
% The course writes these in words on the slides; these render them compactly
% and, more importantly, IDENTICALLY across every problem file.
\newcommand{\CPUtime}{T_{\mathrm{CPU}}}
\newcommand{\IC}{\mathrm{IC}}                    % instruction count
\newcommand{\CPI}{\mathrm{CPI}}
\newcommand{\cycletime}{t_{\mathrm{cycle}}}
\newcommand{\clockrate}{f_{\mathrm{clk}}}
\newcommand{\cycles}{N_{\mathrm{cycles}}}
\newcommand{\speedup}{S}

% Amdahl's law, in the form the slides state it (Module 01, PDF p86, slide 49).
\newcommand{\amdahl}[3]{#1 + \dfrac{#2}{#3}}

%   \begin{document}
%
% Build-check one file, or a whole assignment:
%   docs/latex/build_tex.sh content/cesc_470/hw/hw01/p01_five_components.tex
%   docs/latex/build_tex.sh content/cesc_470/hw/hw01
%
% Doctrine: docs/directives/coursework-solutions.md
% ─────────────────────────────────────────────────────────────────────────────

\documentclass[11pt]{article}

\usepackage[margin=1in]{geometry}
\usepackage{amsmath, amssymb, mathtools}
\usepackage{graphicx}
\usepackage{float}
\usepackage{booktabs}
\usepackage{longtable}
\usepackage{enumitem}
\usepackage{siunitx}
\usepackage[dvipsnames]{xcolor}
\usepackage{tcolorbox}
\usepackage{fancyhdr}
\usepackage{caption}
\usepackage{tikz}
\usepackage{pgfplots}
\pgfplotsset{compat=1.18}
\usepackage[colorlinks=true, allcolors=NavyBlue]{hyperref}

\captionsetup{font=small, labelfont=bf}
\setlist[enumerate]{itemsep=2pt, topsep=4pt}
\setlist[itemize]{itemsep=2pt, topsep=4pt}

% --- final-answer box -------------------------------------------------------
% Every problem file ends with its result in one of these. The condensed
% solutions document is assembled by lifting exactly these.
\newtcolorbox{answerbox}[1][]{
  colback=Green!6, colframe=Green!55!black,
  boxrule=0.6pt, arc=2pt, left=6pt, right=6pt, top=4pt, bottom=4pt,
  #1
}

% --- citation to course material --------------------------------------------
% \cite{Module 01, PDF p55}  ->  a small italic parenthetical.
% Solutions cite the SLIDE/PAGE a formula came from so a grader (or a future
% reader) can check the claim against what was actually taught. See the
% directive: every non-obvious step carries one.
\newcommand{\srcref}[1]{{\small\itshape(#1)}}

% --- a step that came from OUTSIDE the course materials ---------------------
% Used only where the course is genuinely silent and the problem asks for
% outside research. Marked visually so it is never mistaken for lecture content.
\newcommand{\extref}[1]{{\small\itshape\color{Sepia}[external: #1]}}

% --- callout for the trap in a problem --------------------------------------
\newtcolorbox{trap}[1][]{
  colback=BurntOrange!7, colframe=BurntOrange!70!black,
  boxrule=0.6pt, arc=2pt, left=6pt, right=6pt, top=4pt, bottom=4pt,
  fonttitle=\bfseries, title=Watch out, #1
}

% --- per-problem header -----------------------------------------------------
% \problemheader{8}{15 pts}{Target clock rate}
% The optional 4-argument form carries a learning outcome where the course
% states one (CESC 410 does; CESC 470's HW1 does not).
\newcommand{\problemheader}[3]{%
  \begin{center}\large\bfseries
    Problem #1 \quad\normalsize\mdseries(#2)\\[2pt]
    \large\bfseries #3
  \end{center}
  \vspace{2pt}\hrule\vspace{8pt}
}
\newcommand{\problemheaderlo}[4]{%
  \begin{center}\large\bfseries
    Problem #1 \quad\normalsize\mdseries(#2, #3)\\[2pt]
    \large\bfseries #4
  \end{center}
  \vspace{2pt}\hrule\vspace{8pt}
}

% Course name in the running head. Defined BEFORE \lhead references it, and
% \renewcommand'd by each course's macros file (inputted after this one).
\providecommand{\coursename}{}

\pagestyle{fancy}
\fancyhf{}
\lhead{\small\coursename}
\rhead{\small\thepage}
\renewcommand{\headrulewidth}{0.3pt}

%   % CESC 470 Computer Architecture — course-specific macros ONLY.
%
% Inputted AFTER docs/latex/coursework_preamble.tex, which carries everything
% shared (answerbox, problemheader, srcref, page style). Put nothing general
% here: if another course would want it, it belongs in the shared preamble.

\renewcommand{\coursename}{CESC 470}

% --- performance-equation notation ------------------------------------------
% The course writes these in words on the slides; these render them compactly
% and, more importantly, IDENTICALLY across every problem file.
\newcommand{\CPUtime}{T_{\mathrm{CPU}}}
\newcommand{\IC}{\mathrm{IC}}                    % instruction count
\newcommand{\CPI}{\mathrm{CPI}}
\newcommand{\cycletime}{t_{\mathrm{cycle}}}
\newcommand{\clockrate}{f_{\mathrm{clk}}}
\newcommand{\cycles}{N_{\mathrm{cycles}}}
\newcommand{\speedup}{S}

% Amdahl's law, in the form the slides state it (Module 01, PDF p86, slide 49).
\newcommand{\amdahl}[3]{#1 + \dfrac{#2}{#3}}

%   \begin{document}
%
% Build-check one file, or a whole assignment:
%   docs/latex/build_tex.sh content/cesc_470/hw/hw01/p01_five_components.tex
%   docs/latex/build_tex.sh content/cesc_470/hw/hw01
%
% Doctrine: docs/directives/coursework-solutions.md
% ─────────────────────────────────────────────────────────────────────────────

\documentclass[11pt]{article}

\usepackage[margin=1in]{geometry}
\usepackage{amsmath, amssymb, mathtools}
\usepackage{graphicx}
\usepackage{float}
\usepackage{booktabs}
\usepackage{longtable}
\usepackage{enumitem}
\usepackage{siunitx}
\usepackage[dvipsnames]{xcolor}
\usepackage{tcolorbox}
\usepackage{fancyhdr}
\usepackage{caption}
\usepackage{tikz}
\usepackage{pgfplots}
\pgfplotsset{compat=1.18}
\usepackage[colorlinks=true, allcolors=NavyBlue]{hyperref}

\captionsetup{font=small, labelfont=bf}
\setlist[enumerate]{itemsep=2pt, topsep=4pt}
\setlist[itemize]{itemsep=2pt, topsep=4pt}

% --- final-answer box -------------------------------------------------------
% Every problem file ends with its result in one of these. The condensed
% solutions document is assembled by lifting exactly these.
\newtcolorbox{answerbox}[1][]{
  colback=Green!6, colframe=Green!55!black,
  boxrule=0.6pt, arc=2pt, left=6pt, right=6pt, top=4pt, bottom=4pt,
  #1
}

% --- citation to course material --------------------------------------------
% \cite{Module 01, PDF p55}  ->  a small italic parenthetical.
% Solutions cite the SLIDE/PAGE a formula came from so a grader (or a future
% reader) can check the claim against what was actually taught. See the
% directive: every non-obvious step carries one.
\newcommand{\srcref}[1]{{\small\itshape(#1)}}

% --- a step that came from OUTSIDE the course materials ---------------------
% Used only where the course is genuinely silent and the problem asks for
% outside research. Marked visually so it is never mistaken for lecture content.
\newcommand{\extref}[1]{{\small\itshape\color{Sepia}[external: #1]}}

% --- callout for the trap in a problem --------------------------------------
\newtcolorbox{trap}[1][]{
  colback=BurntOrange!7, colframe=BurntOrange!70!black,
  boxrule=0.6pt, arc=2pt, left=6pt, right=6pt, top=4pt, bottom=4pt,
  fonttitle=\bfseries, title=Watch out, #1
}

% --- per-problem header -----------------------------------------------------
% \problemheader{8}{15 pts}{Target clock rate}
% The optional 4-argument form carries a learning outcome where the course
% states one (CESC 410 does; CESC 470's HW1 does not).
\newcommand{\problemheader}[3]{%
  \begin{center}\large\bfseries
    Problem #1 \quad\normalsize\mdseries(#2)\\[2pt]
    \large\bfseries #3
  \end{center}
  \vspace{2pt}\hrule\vspace{8pt}
}
\newcommand{\problemheaderlo}[4]{%
  \begin{center}\large\bfseries
    Problem #1 \quad\normalsize\mdseries(#2, #3)\\[2pt]
    \large\bfseries #4
  \end{center}
  \vspace{2pt}\hrule\vspace{8pt}
}

% Course name in the running head. Defined BEFORE \lhead references it, and
% \renewcommand'd by each course's macros file (inputted after this one).
\providecommand{\coursename}{}

\pagestyle{fancy}
\fancyhf{}
\lhead{\small\coursename}
\rhead{\small\thepage}
\renewcommand{\headrulewidth}{0.3pt}

% CESC 470 Computer Architecture — course-specific macros ONLY.
%
% Inputted AFTER docs/latex/coursework_preamble.tex, which carries everything
% shared (answerbox, problemheader, srcref, page style). Put nothing general
% here: if another course would want it, it belongs in the shared preamble.

\renewcommand{\coursename}{CESC 470}

% --- performance-equation notation ------------------------------------------
% The course writes these in words on the slides; these render them compactly
% and, more importantly, IDENTICALLY across every problem file.
\newcommand{\CPUtime}{T_{\mathrm{CPU}}}
\newcommand{\IC}{\mathrm{IC}}                    % instruction count
\newcommand{\CPI}{\mathrm{CPI}}
\newcommand{\cycletime}{t_{\mathrm{cycle}}}
\newcommand{\clockrate}{f_{\mathrm{clk}}}
\newcommand{\cycles}{N_{\mathrm{cycles}}}
\newcommand{\speedup}{S}

% Amdahl's law, in the form the slides state it (Module 01, PDF p86, slide 49).
\newcommand{\amdahl}[3]{#1 + \dfrac{#2}{#3}}


\begin{document}

\problemheader{3}{5 pts}{How many instructions fit in 32 bits}

\subsection*{Problem statement}

Machine code instructions are stored as bit sequences. If an instruction is 32
bits wide, how many unique instructions can be represented?

\subsection*{Approach}

The premise is the course's own framing of machine code
\srcref{Module 01, PDF p16, slide 7}:

\begin{quote}
``Each machine code instruction consists of a pattern of `1' and `0s' which
determine operation/action to be performed.''
\end{quote}

So the question is purely combinatorial: how many distinct patterns of 32
binary digits exist? Each of the 32 bit positions is an independent choice
between two values, so by the multiplication principle the count is $2$
multiplied by itself 32 times.

\subsection*{Work}

With $b$ bits, each position independently takes one of 2 values:
\begin{equation}
  \underbrace{2 \times 2 \times \cdots \times 2}_{b \text{ factors}} = 2^{b}.
\end{equation}

For $b = 32$:
\begin{equation}
  2^{32} = 4{,}294{,}967{,}296 \approx 4.29 \times 10^{9}.
\end{equation}

A convenient way to evaluate it without a calculator, using $2^{10} = 1024
\approx 10^{3}$:
\begin{equation}
  2^{32} = 2^{2} \times 2^{30} = 4 \times (2^{10})^{3}
         = 4 \times 1024^{3} = 4{,}294{,}967{,}296.
\end{equation}
In the usual binary shorthand this is $4\,\text{Gi}$ — about 4.3 billion.

\begin{trap}
$2^{32}$, not $32^{2}$ and not $2 \times 32$. The exponent is the number of bit
\emph{positions}; the base is the number of values \emph{per} position.
\end{trap}

\paragraph{What this bound actually means in practice.}
$2^{32}$ is the number of distinct 32-bit \emph{patterns} — the absolute ceiling
on encodable instructions. A real ISA never gets close, because the 32 bits are
partitioned into fields: an opcode, register specifiers, and immediate values
\srcref{PDF p23, slide 11 — ``Encoding and representing instructions''}. For
example, if 6 bits are spent on the opcode, only $2^{6} = 64$ distinct
operations are encodable, and the remaining 26 bits carry operands. So the
answer to the question as asked is $2^{32}$ \emph{bit patterns}, which is an
upper bound on, not an estimate of, an ISA's instruction count.

\subsection*{Check}

Dimensional sanity: 32 bits is 4 bytes, and 4 bytes is exactly the size of an
unsigned 32-bit integer, whose range is $0$ to $2^{32}-1$ — that is $2^{32}$
distinct values, matching. Independently computed: $2^{32} = 4{,}294{,}967{,}296$.

\subsection*{Answer}

\begin{answerbox}
$$2^{32} = 4{,}294{,}967{,}296 \approx 4.29\times10^{9} \text{ unique instructions.}$$
Each of the 32 bit positions independently takes 2 values, giving $2^{32}$
distinct patterns. This is an \emph{upper bound}: a real ISA divides the 32 bits
into opcode and operand fields, so far fewer distinct \emph{operations} are
actually encoded.
\end{answerbox}

\end{document}
